% Propuesta — Proyecto 2 · Visualización I · Equipo 7
% Estilo académico (artículo) — sin colores ni decoración
% Compilar: pdflatex Propuesta_Equipo7.tex (2 pasadas)
\documentclass[11pt,a4paper]{article}

\usepackage[utf8]{inputenc}
\usepackage[T1]{fontenc}
\usepackage[spanish]{babel}
\usepackage{lmodern}
\usepackage[margin=2.5cm]{geometry}
\usepackage{amsmath,amssymb}
\usepackage{booktabs}
\usepackage{array}
\usepackage{enumitem}
\usepackage{caption}
\usepackage[hidelinks]{hyperref}

% --- Estructura académica: secciones numeradas en negrita, sin color ---
\setlength{\parindent}{0pt}
\setlength{\parskip}{0.55em}
\setcounter{secnumdepth}{2}
\makeatletter
\renewcommand\section{\@startsection{section}{1}{\z@}%
  {-1.8ex \@plus -0.5ex \@minus -0.2ex}%
  {0.9ex \@plus0.2ex}%
  {\normalfont\large\bfseries}}
\renewcommand\subsection{\@startsection{subsection}{2}{\z@}%
  {-1.5ex \@plus -0.4ex \@minus -0.2ex}%
  {0.7ex \@plus0.2ex}%
  {\normalfont\normalsize\bfseries}}
\makeatother

% Tablas académicas (sin líneas verticales, booktabs)
\captionsetup{font=small,labelfont=bf}
\renewcommand{\arraystretch}{1.2}

\begin{document}

% ---------------- Portada académica ----------------
\begin{center}
{\LARGE\bfseries Propuesta de investigación --- Proyecto 2}\\[1.2em]
{\large Visualización reproducible en Python}\\[0.6em]
{\large Curso de Visualización I · Quinto Semestre de 2026}\\[1.6em]
{\large Joshua Trujillo y Nicolás Figueroa}\\[0.3em]
{\large (Equipo~7)}\\[1.6em]
{\normalsize Docentes: Adriana Pineda y Paola Moreno}\\[0.3em]
{\normalsize Universidad Externado de Colombia}\\[0.3em]
{\normalsize 12 de octubre de 2026}
\end{center}

\vspace{1.5em}
\noindent\textbf{Entrega máxima del proyecto:} 26 de octubre de 2026.\\
\textbf{Estado de la propuesta:} enviada para aprobación docente.
\vspace{0.8em}

\hrule
\vspace{1em}

% ---------------- 1 ----------------
\section{Audiencia técnica específica}

\textbf{Destinatario.} Unidad de investigación comparativa sobre determinantes del
bienestar de una organización internacional (p.\ ej., el grupo de investigación
sobre desarrollo humano y calidad de vida del PNUD o del Banco Mundial).

\textbf{Interés en la problemática.} El equipo requiere identificar, de manera
fundamentada y reproducible, qué combinaciones de condiciones ---PIB per cápita,
apoyo social, esperanza de vida saludable, libertad, generosidad y percepción de
corrupción--- se asocian con niveles altos de bienestar reportado, y en qué medida
dichas condiciones no alcanzan a explicar la evaluación de vida de algunos países.
Los resultados orientan investigaciones posteriores sobre los determinantes
explicados y no explicados del bienestar y ofrecen una base verificable para
estudios comparados.

% ---------------- 2 ----------------
\section{Pregunta analítica}

\begin{quote}
\textbf{¿Qué «recetas» de bienestar predominan entre los países con mayor y menor
evaluación de vida según el \emph{World Happiness Report} 2024, y en qué medida el
componente no explicado (Distopía + residual) altera o cuestiona dichas recetas?}
\end{quote}

\textbf{Subpreguntas.}
\begin{enumerate}[leftmargin=1.6em, itemsep=0.25em]
  \item ¿Qué combinación de los seis indicadores $c_{i,k}$ caracteriza a los
        países de los terciles superior e inferior del \emph{ladder score} $L_i$?
  \item ¿Qué países presentan el componente $D_i$ (\texttt{Dystopia + residual})
        más alto, y qué tienen en común, sugiriendo que las variables disponibles
        no bastan para explicar su bienestar?
  \item ¿Los países atípicos respecto de su «receta» esperada concentran valores
        elevados o reducidos del residual, lo que indicaría una explicación
        incompleta?
\end{enumerate}

% ---------------- 3 ----------------
\section{Definición de «receta de bienestar»}

\textbf{Definición operativa.} Una \emph{«receta de bienestar»} es el perfil
cuantitativo de contribuciones con que un país alcanza su evaluación de vida: la
combinación relativa de los seis indicadores explicativos del \emph{World Happiness
Report}, esto es, cuánto aporta cada dimensión al \emph{ladder score} nacional.
Dos países pueden alcanzar puntuaciones similares con recetas distintas ---por
ejemplo, uno sustentado en PIB y esperanza de vida, otro en apoyo social y
libertad---. Estudiar las recetas implica comparar qué dimensiones predominan en
los países de mayor y menor bienestar, y no únicamente sus puntuaciones absolutas.

\textbf{Formalización.} Sea $L_i$ el \emph{ladder score} del país $i$. Su receta
es el vector de contribuciones
\begin{equation}
  R_i \;=\; \bigl(\, c_{i,\text{PIB}},\; c_{i,\text{social}},\;
  c_{i,\text{vida}},\; c_{i,\text{libertad}},\; c_{i,\text{generosidad}},\;
  c_{i,\text{corrupción}} \,\bigr),
\end{equation}
de modo que se cumple la descomposición aditiva
\begin{equation}
  L_i \;=\; \sum_{k=1}^{6} c_{i,k} \;+\; D_i,
  \label{eq:descomposicion}
\end{equation}
donde $D_i$ denota el componente no explicado por los seis indicadores. La
composición relativa $c_{i,k}\big/\sum_{k} c_{i,k}$ permite comparar recetas entre
países con distintos niveles de $L_i$.

\textbf{Advertencia inferencial.} Las contribuciones $c_{i,k}$ son componentes
aditivos de una descomposición estadística del \emph{World Happiness Report}; las
asociaciones observadas entre dichos componentes y $L_i$ no implican causalidad.

% ---------------- 4 ----------------
\section{Definición de las variables}

Cada variable se define tal como aparece en el archivo oficial
\texttt{WHR24\_Data\_Figure\_2.1.xls} (\emph{Data for Figure 2.1}, \emph{World
Happiness Report} 2024). Se emplean las notaciones $L_i$, $c_{i,k}$ y $D_i$
introducidas en la sección~3.

\subsection{Variable de interés}

\begin{description}[style=nextline, leftmargin=1.7em, itemsep=0.35em]
  \item[\texttt{Ladder score} ($L_i$)] Evaluación de vida promedio reportada por
    los habitantes de cada país en la escala Cantril del \emph{World Happiness
    Report}: valores de 0 (peor vida posible) a 10 (mejor vida posible), a partir
    de las encuestas Gallup 2021--2023. Es la variable a la que la pregunta alude
    como «evaluación de vida» y el criterio de orden para definir los países «con
    mayor y menor bienestar». Rango observado en 2024: 1{,}72 (Afganistán) a
    7{,}74 (Finlandia).
\end{description}

\subsection{Los seis indicadores de la receta}

Cada columna \texttt{Explained by: X} corresponde a la \textbf{contribución} del
factor respectivo al \emph{ladder score} del país, expresada en puntos de la misma
escala 0--10; la suma de los seis componentes más el residual reconstruye $L_i$
(ec.~\ref{eq:descomposicion}). No se trata de porcentajes ni de puntajes
independientes.

\begin{description}[style=nextline, leftmargin=2.6em, itemsep=0.35em]
  \item[\texttt{Log GDP per capita} ($c_{i,\text{PIB}}$)]
    Contribución del logaritmo natural del PIB per cápita en paridad de poder de
    compra. Dimensión~1 de la receta.
  \item[\texttt{Social support} ($c_{i,\text{social}}$)]
    Contribución del apoyo social: proporción de personas con familiares o amigos
    de confianza a quienes acudir en momentos de necesidad. Dimensión~2.
  \item[\texttt{Healthy life expectancy} ($c_{i,\text{vida}}$)]
    Contribución de la esperanza de vida saludable al nacer, en años ajustados por
    discapacidad (GBD). Dimensión~3.
  \item[\texttt{Freedom to make life choices} ($c_{i,\text{libertad}}$)]
    Contribución de la libertad percibida para tomar decisiones de vida.
    Dimensión~4.
  \item[\texttt{Generosity} ($c_{i,\text{generosidad}}$)]
    Contribución de la generosidad: correlación ajustada por ingreso entre
    donaciones benéficas recientes (último mes) y PIB per cápita. Dimensión~5.
  \item[\texttt{Perceptions of corruption} ($c_{i,\text{corrupción}}$)]
    Contribución de la percepción de corrupción (gobierno y empresas): un valor
    más alto indica menor percepción de corrupción y, por tanto, mayor
    contribución. Dimensión~6.
\end{description}

\subsection{Componente no explicado}

\begin{description}[style=nextline, leftmargin=2.6em, itemsep=0.35em]
  \item[\texttt{Dystopia + residual} ($D_i$)]
    Suma de la puntuación de \emph{Distopía} ---país hipotético con los valores
    mínimos observados de los seis factores, usado como referencia--- y del
    residual no explicado: la parte del \emph{ladder score} que los seis
    indicadores no alcanzan a cubrir. Constituye el eje de la subpregunta~2:
    valores elevados señalan países cuyo bienestar queda fuera de la explicación
    estándar.
\end{description}

\subsection{Variables de contexto}

\begin{description}[style=nextline, leftmargin=2.6em, itemsep=0.35em]
  \item[\texttt{Country name}]
    Nombre del país ($n = 143$ observaciones). Se emplea para identificación y
    anotaciones en los gráficos; no es una variable numérica.
  \item[\texttt{upperwhisker}, \texttt{lowerwhisker}]
    Extremos del intervalo de confianza al 95\% de $L_i$. Se usan como bandas de
    incertidumbre de referencia; no integran la receta.
\end{description}

\textbf{Valores faltantes.} Tres países (Bahrain, Tajikistan y State of Palestine)
presentan faltantes en los seis indicadores y en $D_i$, aunque sí disponen de
$L_i$. Se conservan en el conjunto de datos y se excluyen de manera explícita
únicamente en las visualizaciones que requieran valores completos, documentando
dicho filtro.

% ---------------- 5 ----------------
\section{Tabla resumen de variables}

\begin{table}[htbp]
\centering
\caption{Variables del archivo oficial y su uso en el análisis.}
\label{tab:variables}
\small
\begin{tabular}{>{\raggedright\arraybackslash}p{4.9cm}
                >{\raggedright\arraybackslash}p{4.7cm}
                >{\raggedright\arraybackslash}p{5.1cm}}
\toprule
\textbf{Variable (archivo oficial)} & \textbf{Significado} & \textbf{Uso previsto} \\
\midrule
\texttt{Country name} & País & Identificación y anotaciones \\
\texttt{Ladder score} ($L_i$) & Evaluación de vida (0--10) & Variable de interés; eje de orden \\
\texttt{upperwhisker}, \texttt{lowerwhisker} & Intervalo de confianza de $L_i$ & Bandas de incertidumbre \\
\texttt{Explained by: Log GDP per capita} & Contribución del PIB per cápita & Dimensión~1 de la receta \\
\texttt{Explained by: Social support} & Contribución del apoyo social & Dimensión~2 \\
\texttt{Explained by: Healthy life expectancy} & Contribución de la esperanza de vida saludable & Dimensión~3 \\
\texttt{Explained by: Freedom to make life choices} & Contribución de la libertad & Dimensión~4 \\
\texttt{Explained by: Generosity} & Contribución de la generosidad & Dimensión~5 \\
\texttt{Explained by: Perceptions of corruption} & Contribución de la percepción de corrupción & Dimensión~6 \\
\texttt{Dystopia + residual} ($D_i$) & Componente no explicado & Eje de la subpregunta~2 \\
\bottomrule
\end{tabular}
\end{table}

\textbf{Criterio de selección.} Se emplean las variables numéricas disponibles en
el archivo oficial. No se construirán indicadores derivados, salvo la
clasificación en terciles del \emph{ladder score} (\texttt{pandas.qcut},
\texttt{q = 3}), cuyo criterio se documentará en el notebook.

% ---------------- 6 ----------------
\section{Plan de visualizaciones}

\begin{table}[htbp]
\centering
\caption{Anteproyecto de visualizaciones: cuatro estáticas (Seaborn) y dos
interactivas (Plotly), conforme a los requisitos del enunciado.}
\label{tab:plan}
\small
\begin{tabular}{>{\raggedright\arraybackslash}p{2.0cm}
                >{\raggedright\arraybackslash}p{7.5cm}
                >{\raggedright\arraybackslash}p{5.2cm}}
\toprule
\textbf{\#} & \textbf{Visualización} & \textbf{Propósito analítico} \\
\midrule
Seaborn 1 & Distribución de $L_i$ y de $D_i$ & Forma global y extremos \\
Seaborn 2 & Diagramas de caja de los seis indicadores por tercil de bienestar & Comparación de recetas \\
Seaborn 3 & Mapa de calor de la composición relativa por grupo & Perfil de receta \\
Seaborn 4 & $D_i$ frente a $L_i$ & Componente no explicado \\
Plotly 1 & Diagrama de dispersión interactivo (tamaño = $D_i$, anotación de países) & Exploración de atípicos \\
Plotly 2 & Barras \emph{top-15} con selector de indicador y color por tercil & Exploración por dimensión \\
\bottomrule
\end{tabular}
\end{table}

% ---------------- 7 ----------------
\section{Coherencia del análisis}

La secuencia de trabajo conserva la coherencia exigida por el enunciado:
audiencia (unidad de investigación comparativa) $\rightarrow$ pregunta analítica
(recetas y componente no explicado) $\rightarrow$ variables (seis indicadores,
residual y \emph{ladder score}) $\rightarrow$ preparación (renombre documentado,
terciles, manejo de faltantes) $\rightarrow$ visualizaciones (composición por
grupos, perfiles, residual y atípicos) $\rightarrow$ interpretación escrita tras
cada visualización $\rightarrow$ conclusiones sustentadas en la evidencia.

\end{document}
